\documentclass[10pt,a4paper]{amsart}
\usepackage[margin=19mm]{geometry}
\usepackage{amsmath,amssymb,mathtools}
\usepackage[scr=rsfs,cal=euler]{mathalfa}
\usepackage{xfrac}
\usepackage[unicode,hidelinks,bookmarks=false]{hyperref}
\newcommand{\ie}{i.e.\ }
\newcommand{\cf}{cf.\ }
\newcommand{\resp}{respectively\ }
\newcommand{\Subsection}{\subsection}
\providecommand{\nobreakdash}{\nobreak\hbox{-}}
\setlength{\emergencystretch}{2em}
\multlinegap0pt
\usepackage{float}
\usepackage{bm}
\usepackage{amssymb}
\usepackage{enumerate}
\usepackage{xcolor}
\newtheorem{assumption}{Assumption}
\newtheorem{proposition}{Proposition}
\newcommand{\C}{\mathbb{C}}
\newcommand{\R}{\mathbb{R}}
\newcommand{\ii}{\mathrm{i}}
\title{Rough Heston model as the scaling limit of bivariate cumulative heavy-tailed INAR processes: Weak-error bounds and option pricing}
\author{Yingli Wang, Zhenyu Cui, Lingjiong Zhu}
\date{Source: arXiv:2503.18259v6 (CC BY 4.0)}
\begin{document}
\maketitle

\noindent\emph{Source and licence.} Rough Heston model as the scaling limit of bivariate cumulative heavy-tailed INAR processes: Weak-error bounds and option pricing. Version arXiv:2503.18259v6 (CC BY 4.0), distributed by arXiv under the Creative Commons Attribution 4.0 International licence, http://creativecommons.org/licenses/by/4.0/, as recorded in the arXiv licence metadata for this exact version and shown on the abstract page https://arxiv.org/abs/2503.18259v6. Full licence notice: \texttt{licenses/arxiv-2503.18259-CC-BY-4.0.txt}.\par\smallskip

\noindent\textit{Editorial scope.} Counted unit: the article's proposition prop:riccati-regularity (source0.tex lines 5048--5068). The article's complete original proof of it is delivered with it (source0.tex lines 5070--5575, one \texttt{proof} environment of 506 lines). No mathematics is written, repaired or re-derived: the statement and the proof are the article's own lines, in the article's own notation, and no retained line is substituted or re-flowed.  Retained uncounted as the source-local context the proof needs: lines 724--732.  Nothing was omitted from any retained span, so the package carries no omission record.  No retained line contains a citation, so the wrapper carries no bibliography and no bibliography entry.  No retained line contains a figure environment, an image-inclusion command or drawing code, so no image is delivered and the package declares no asset dependency.  Editorial formatting only, in addition: an AMS \texttt{amsart} preamble that restates the article's own declarations and macros used by the retained lines, namely the environments \texttt{assumption}, \texttt{proposition} on separate counters, together with the standard packages those lines need.  The retained environments print in this document as Assumption 1, Proposition 1 in order of appearance, whereas the article prints the same items with the numbering of its own full document.  The complete native source of the article as distributed in the arXiv e-print for this exact version is bundled unchanged as \texttt{sources/175246/source0.tex}.
\begin{assumption}[Bounded-strip Riccati regularity]\label{ass:riccati-bound}
Fix $\eta\in\R$ and $R>0$ with $\Gamma_{\eta,R}\subset\mathcal S_T$. We assume that the continuous fractional Riccati solution satisfies the bounded-strip supremum bound
\[
  M_{\eta,R,T}
  :=
  \sup_{z\in\Gamma_{\eta,R}}\sup_{0\le t\le T}|h(t,z)|
  <\infty.
\]
\end{assumption}

\begin{proposition}[Bounded-strip regularity of the continuous fractional Riccati solution]
\label{prop:riccati-regularity}
Fix $\eta\in\R$ and $R>0$ such that
\[
  \Gamma_{\eta,R}:=\{\eta+\ii\xi:\ |\xi|\le R\}\subset \mathcal S_T.
\]
Under Assumption~\ref{ass:riccati-bound}, the map $(t,z)\mapsto h(t,z)$ is
jointly continuous on $[0,T]\times\Gamma_{\eta,R}$, each map
$t\mapsto h(t,z)$ belongs to $AC([0,T])\cap C^1((0,T])$, and there exists a
constant $C_{\eta,R,T}>0$ such that, uniformly for $z\in\Gamma_{\eta,R}$,
\begin{align}
  |h(t,z)| &\le C_{\eta,R,T}t^\alpha,
  &&0\le t\le T,\label{eq:h-small-t-prop}\\
  |h(t,z)-h(s,z)| &\le C_{\eta,R,T}|t-s|^\alpha,
  &&0\le s,t\le T,\label{eq:h-holder-t-prop}\\
  |\partial_t h(t,z)| &\le C_{\eta,R,T}t^{\alpha-1},
  &&0<t\le T.\label{eq:h-derivative-t-prop}
\end{align}
Moreover, the map $z\mapsto h(\cdot,z)$ is locally Lipschitz from
$\Gamma_{\eta,R}$ into $C([0,T])$.
\end{proposition}

\begin{proof}
Throughout the proof, $C_{\eta,R,T}>0$ denotes a generic constant that may
change from line to line but depends only on $(\eta,R,T)$ and the fixed model
parameters.

Let
\[
  F_z(x):=\frac12(z^2-z)+\gamma(\rho\nu z-1)x+\frac{(\gamma\nu)^2}{2}x^2,
  \qquad x\in\C.
\]
Then the Volterra form of the fractional Riccati equation reads
\begin{equation}\label{eq:riccati-volterra-proof-reg}
  h(t,z)
  =
  \frac{1}{\Gamma(\alpha)}
  \int_0^t (t-s)^{\alpha-1}F_z(h(s,z))\,ds,
  \qquad 0\le t\le T.
\end{equation}

Since $\Gamma_{\eta,R}$ is compact and
\[
  \sup_{z\in\Gamma_{\eta,R}}\sup_{0\le t\le T}|h(t,z)|
  \le M_{\eta,R,T},
\]
the polynomial $F_z(x)$ and its derivative in $x$ are uniformly bounded on the
set
\[
  \{(z,x): z\in\Gamma_{\eta,R},\ |x|\le M_{\eta,R,T}\}.
\]
Hence there exists a constant $C_{\eta,R,T}>0$ such that
\begin{equation}\label{eq:F-uniform-bound-proof-reg}
  |F_z(h(t,z))|\le C_{\eta,R,T},
  \qquad
  0\le t\le T,\ z\in\Gamma_{\eta,R},
\end{equation}
and
\begin{equation}\label{eq:F-lipschitz-proof-reg}
  |F_z(x)-F_z(y)|\le C_{\eta,R,T}|x-y|,
\end{equation}
for all $z\in\Gamma_{\eta,R}$ and all $x,y$ with $|x|,|y|\le M_{\eta,R,T}$.

We divide the proof into four steps.

\medskip
\noindent
\textit{Step 1: small-time bound.}
By \eqref{eq:riccati-volterra-proof-reg} and \eqref{eq:F-uniform-bound-proof-reg},
\begin{align*}
  |h(t,z)|
  &\le
  \frac{1}{\Gamma(\alpha)}
  \int_0^t (t-s)^{\alpha-1}|F_z(h(s,z))|\,ds\\
  &\le
  \frac{C_{\eta,R,T}}{\Gamma(\alpha)}
  \int_0^t (t-s)^{\alpha-1}\,ds
  =
  \frac{C_{\eta,R,T}}{\Gamma(\alpha+1)}\,t^\alpha.
\end{align*}
This proves \eqref{eq:h-small-t-prop}.

\medskip
\noindent
\textit{Step 2: $\alpha$-H\"older continuity in time.}
Fix $0\le s<t\le T$. Subtracting the two Volterra formulas gives
\begin{align*}
  h(t,z)-h(s,z)
  &=
  \frac{1}{\Gamma(\alpha)}
  \int_0^s
  \left((t-u)^{\alpha-1}-(s-u)^{\alpha-1}\right)F_z(h(u,z))\,du\\
  &\quad
  +
  \frac{1}{\Gamma(\alpha)}
  \int_s^t (t-u)^{\alpha-1}F_z(h(u,z))\,du.
\end{align*}
Taking absolute values and using \eqref{eq:F-uniform-bound-proof-reg},
\begin{align*}
  |h(t,z)-h(s,z)|
  &\le
  \frac{C_{\eta,R,T}}{\Gamma(\alpha)}
  \int_0^s
  \left|(t-u)^{\alpha-1}-(s-u)^{\alpha-1}\right|du
  \\
  &\qquad\qquad\qquad
  +
  \frac{C_{\eta,R,T}}{\Gamma(\alpha)}
  \int_s^t (t-u)^{\alpha-1}du.
\end{align*}
Since $\alpha-1\in(-1,0)$, the map $x\mapsto x^{\alpha-1}$ is decreasing on
$(0,\infty)$, and therefore
\begin{align*}
  \int_0^s\left((s-u)^{\alpha-1}-(t-u)^{\alpha-1}\right)\,du
  &=
  \frac{1}{\alpha}\left(s^\alpha-t^\alpha+(t-s)^\alpha\right)
  \le \frac{1}{\alpha}(t-s)^\alpha.
\end{align*}
Also,
$\int_s^t (t-u)^{\alpha-1}\,du=\frac1\alpha (t-s)^\alpha$.
Hence,
\[
  |h(t,z)-h(s,z)|
  \le C_{\eta,R,T}|t-s|^\alpha,
\]
uniformly in $z\in\Gamma_{\eta,R}$. This proves \eqref{eq:h-holder-t-prop}.

\medskip
\noindent
\textit{Step 3: continuity and Lipschitz dependence on $z$.}
Fix $z_1,z_2\in\Gamma_{\eta,R}$ and write
\[
  \delta h(t):=h(t,z_1)-h(t,z_2).
\]
Subtracting the two Volterra equations gives
\begin{align*}
  \delta h(t)
  &=
  \frac{1}{\Gamma(\alpha)}
  \int_0^t (t-s)^{\alpha-1}
  \left(F_{z_1}(h(s,z_1))-F_{z_2}(h(s,z_2))\right)\,ds.
\end{align*}
Split the integrand as
\begin{align*}
  F_{z_1}(h(s,z_1))-F_{z_2}(h(s,z_2))
  &=
  \left(F_{z_1}(h(s,z_1))-F_{z_1}(h(s,z_2))\right)
  \\
  &\quad+
  \left(F_{z_1}(h(s,z_2))-F_{z_2}(h(s,z_2))\right).
\end{align*}
By \eqref{eq:F-lipschitz-proof-reg},
\[
  \left|F_{z_1}(h(s,z_1))-F_{z_1}(h(s,z_2))\right|
  \le C_{\eta,R,T}|\delta h(s)|.
\]
Since $F_z(x)$ is polynomial in $(z,x)$ and $z\in\Gamma_{\eta,R}$, $|x|\le
M_{\eta,R,T}$, there exists $C_{\eta,R,T}>0$ such that
\[
  \left|F_{z_1}(h(s,z_2))-F_{z_2}(h(s,z_2))\right|
  \le C_{\eta,R,T}|z_1-z_2|.
\]
Therefore,
\[
  |\delta h(t)|
  \le
  C_{\eta,R,T}|z_1-z_2|t^\alpha
  +
  C_{\eta,R,T}
  \int_0^t (t-s)^{\alpha-1}|\delta h(s)|\,ds.
\]
Set
\[
u(t):=\sup_{0\le s\le t}|h(s,z_1)-h(s,z_2)|.
\]
Then the previous estimate implies
\[
u(t)\le C_{\eta,R,T}|z_1-z_2|
   +C_{\eta,R,T}\int_0^t (t-s)^{\alpha-1}u(s)\,ds,
   \qquad 0\le t\le T.
\]
By the weakly singular Gr\"onwall inequality, it follows that
\[
u(t)\le C_{\eta,R,T}|z_1-z_2|\,E_\alpha(C_{\eta,R,T}t^\alpha),
\qquad 0\le t\le T.
\]
Since $t\in[0,T]$, we obtain
\[
\sup_{0\le t\le T}|h(t,z_1)-h(t,z_2)|
\le C_{\eta,R,T}|z_1-z_2|.
\]
In particular, $z\mapsto h(\cdot,z)$ is locally Lipschitz into $C([0,T])$.
Combining this with \eqref{eq:h-holder-t-prop} gives joint continuity of
$(t,z)\mapsto h(t,z)$ on $[0,T]\times\Gamma_{\eta,R}$.

\medskip
\noindent
\textit{Step 4: absolute continuity and derivative bound.}
Define
\[
  g_z(t):=F_z(h(t,z)),\qquad 0\le t\le T.
\]
By \eqref{eq:h-holder-t-prop} and the local Lipschitz property
\eqref{eq:F-lipschitz-proof-reg},
\[
  |g_z(t)-g_z(s)|\le C_{\eta,R,T}|t-s|^\alpha,
  \qquad 0\le s,t\le T.
\]
Thus $g_z\in C^\alpha([0,T])$ uniformly in $z$, and
$h(\cdot,z)=I^\alpha g_z$.
We now prove that $h(\cdot,z)$ is differentiable on $(0,T]$.
Fix $z\in\Gamma_{\eta,R}$ and $t\in(0,T]$. Recall that
\[
  h(t,z)=\frac{1}{\Gamma(\alpha)}\int_0^t (t-s)^{\alpha-1}g_z(s)\,ds,
\]
where
\[
  |g_z(s)-g_z(r)|\le C_{\eta,R,T}|s-r|^\alpha,
  \qquad 0\le r,s\le T.
\]

Let $h_0$ be such that $|h_0|<\min\{t,T-t\}$, and let
$u\in(-h_0,h_0)\setminus\{0\}$.

\smallskip
\noindent
\emph{Case 1: $u>0$.}
In this case,
\begin{align*}
  \frac{h(t+u,z)-h(t,z)}{u}
  &=
  \frac{1}{\Gamma(\alpha)}
  \int_0^t
  \frac{(t+u-s)^{\alpha-1}-(t-s)^{\alpha-1}}{u}\,g_z(s)\,ds
  \\
  &\quad+
  \frac{1}{\Gamma(\alpha)}
  \frac{1}{u}\int_t^{t+u}(t+u-s)^{\alpha-1}g_z(s)\,ds.
\end{align*}
Adding and subtracting $g_z(t)$, we rewrite this as
\begin{align}
  &\frac{h(t+u,z)-h(t,z)}{u}\nonumber\\
  &=
  \frac{1}{\Gamma(\alpha)}
  \int_0^t
  \frac{(t+u-s)^{\alpha-1}-(t-s)^{\alpha-1}}{u}
  \left(g_z(s)-g_z(t)\right)\,ds
  \nonumber\\
  &\qquad+
  \frac{1}{\Gamma(\alpha)}
  \frac{1}{u}\int_t^{t+u}(t+u-s)^{\alpha-1}
  \left(g_z(s)-g_z(t)\right)\,ds
  \nonumber\\
  &\qquad\qquad+
  \frac{g_z(t)}{\Gamma(\alpha)}
  \left[
    \int_0^t
    \frac{(t+u-s)^{\alpha-1}-(t-s)^{\alpha-1}}{u}\,ds
    +
    \frac{1}{u}\int_t^{t+u}(t+u-s)^{\alpha-1}\,ds
  \right].
  \label{eq:dq-split-proof-pos}
\end{align}

For the bracketed term, a direct computation gives
\begin{align*}
  &\int_0^t
    \frac{(t+u-s)^{\alpha-1}-(t-s)^{\alpha-1}}{u}\,ds
    +
    \frac{1}{u}\int_t^{t+u}(t+u-s)^{\alpha-1}\,ds
  \\
  &=
  \frac{1}{u}\int_0^{t+u}r^{\alpha-1}\,dr
  -
  \frac{1}{u}\int_0^t r^{\alpha-1}\,dr
  =
  \frac{(t+u)^\alpha-t^\alpha}{\alpha u},
\end{align*}
and therefore
\begin{equation}\label{eq:bracket-limit-pos}
  \frac{g_z(t)}{\Gamma(\alpha)}
  \left[
    \int_0^t
    \frac{(t+u-s)^{\alpha-1}-(t-s)^{\alpha-1}}{u}\,ds
    +
    \frac{1}{u}\int_t^{t+u}(t+u-s)^{\alpha-1}\,ds
  \right]
  \to
  \frac{g_z(t)}{\Gamma(\alpha)}\,t^{\alpha-1}
\end{equation}
as $u\downarrow0$.

For the first integral in \eqref{eq:dq-split-proof-pos}, fix $0<u<t/2$. By the mean value theorem,
for each $s\in[0,t)$ there exists $\theta=\theta(s,u)\in(0,1)$ such that
\[
  \frac{(t+u-s)^{\alpha-1}-(t-s)^{\alpha-1}}{u}
  =
  (\alpha-1)(t-s+\theta u)^{\alpha-2}.
\]
Hence
\[
  \left|
    \frac{(t+u-s)^{\alpha-1}-(t-s)^{\alpha-1}}{u}
  \right|
  \le C_\alpha (t-s)^{\alpha-2},
\]
and thus, using the $\alpha$-H\"older continuity of $g_z$,
\[
  \left|
    \frac{(t+u-s)^{\alpha-1}-(t-s)^{\alpha-1}}{u}
    \left(g_z(s)-g_z(t)\right)
  \right|
  \le C_{\eta,R,T}(t-s)^{2\alpha-2}.
\]
Since $\alpha>1/2$, we have $2\alpha-2>-1$, so $(t-s)^{2\alpha-2}\in L^1(0,t)$.
Moreover, for each fixed $s\in[0,t)$,
\[
  \frac{(t+u-s)^{\alpha-1}-(t-s)^{\alpha-1}}{u}
  \to
  (\alpha-1)(t-s)^{\alpha-2}
\]
as $u\downarrow0$. Hence, by the dominated convergence theorem,
\begin{align}
  &\int_0^t
  \frac{(t+u-s)^{\alpha-1}-(t-s)^{\alpha-1}}{u}
  \left(g_z(s)-g_z(t)\right)\,ds
  \nonumber\\
  &\to
  (\alpha-1)\int_0^t (t-s)^{\alpha-2}\left(g_z(s)-g_z(t)\right)\,ds
  \qquad\text{as }u\downarrow0.
  \label{eq:first-integral-limit-pos}
\end{align}

For the second integral in \eqref{eq:dq-split-proof-pos}, again by the $\alpha$-H\"older continuity of $g_z$,
\begin{align}
  &\left|
    \frac{1}{u}\int_t^{t+u}(t+u-s)^{\alpha-1}
    \left(g_z(s)-g_z(t)\right)\,ds
  \right|
  \nonumber\\
  &\le
  \frac{C_{\eta,R,T}}{u}\int_t^{t+u}(t+u-s)^{\alpha-1}(s-t)^\alpha\,ds
  \nonumber\\
 &=
  \frac{C_{\eta,R,T}}{u}\int_0^u r^{\alpha-1}(u-r)^\alpha\,dr
 =
  C_{\eta,R,T}u^{2\alpha-1}\int_0^1 x^{\alpha-1}(1-x)^\alpha\,dx
  \to0
\label{eq:second-integral-limit-pos}
\end{align}
as $u\downarrow0$.

Combining the three limits above in \eqref{eq:bracket-limit-pos}, \eqref{eq:first-integral-limit-pos}, and \eqref{eq:second-integral-limit-pos},
we conclude that
\begin{equation}
  \lim_{u\downarrow0}\frac{h(t+u,z)-h(t,z)}{u}
  =
  \frac{g_z(t)}{\Gamma(\alpha)}t^{\alpha-1}
  +
  \frac{\alpha-1}{\Gamma(\alpha)}
  \int_0^t (t-s)^{\alpha-2}\left(g_z(s)-g_z(t)\right)\,ds .
  \label{eq:right-derivative-proof}
\end{equation}

\smallskip
\noindent
\emph{Case 2: $u<0$.}
In this case,
\begin{align*}
  \frac{h(t+u,z)-h(t,z)}{u}
  &=
  \frac{1}{\Gamma(\alpha)}
  \int_0^{t+u}
  \frac{(t+u-s)^{\alpha-1}-(t-s)^{\alpha-1}}{u}\,g_z(s)\,ds
  \\
  &\quad
  -
  \frac{1}{\Gamma(\alpha)}
  \frac{1}{u}\int_{t+u}^{t}(t-s)^{\alpha-1}g_z(s)\,ds.
\end{align*}
Adding and subtracting $g_z(t)$, we obtain
\begin{align}
  &\frac{h(t+u,z)-h(t,z)}{u}\nonumber\\
  &=
  \frac{1}{\Gamma(\alpha)}
  \int_0^{t+u}
  \frac{(t+u-s)^{\alpha-1}-(t-s)^{\alpha-1}}{u}
  \left(g_z(s)-g_z(t)\right)\,ds
  \nonumber\\
  &\qquad
  -
  \frac{1}{\Gamma(\alpha)}
  \frac{1}{u}\int_{t+u}^{t}(t-s)^{\alpha-1}
  \left(g_z(s)-g_z(t)\right)\,ds
  \nonumber\\
  &\qquad\qquad+
  \frac{g_z(t)}{\Gamma(\alpha)}
  \left[
    \int_0^{t+u}
    \frac{(t+u-s)^{\alpha-1}-(t-s)^{\alpha-1}}{u}\,ds
    -
    \frac{1}{u}\int_{t+u}^{t}(t-s)^{\alpha-1}\,ds
  \right].
  \label{eq:dq-split-proof-neg}
\end{align}

For the bracketed term, another direct computation gives
\begin{align*}
  &\int_0^{t+u}
  \frac{(t+u-s)^{\alpha-1}-(t-s)^{\alpha-1}}{u}\,ds
  -
  \frac{1}{u}\int_{t+u}^{t}(t-s)^{\alpha-1}\,ds
  \\
  &=
  \frac{1}{u}\int_0^{t+u}r^{\alpha-1}\,dr
  -
  \frac{1}{u}\int_0^t r^{\alpha-1}\,dr
  =
  \frac{(t+u)^\alpha-t^\alpha}{\alpha u},
\end{align*}
so again
\begin{equation}\label{eq:bracket-limit-neg}
  \frac{g_z(t)}{\Gamma(\alpha)}
  \left[
    \int_0^{t+u}
    \frac{(t+u-s)^{\alpha-1}-(t-s)^{\alpha-1}}{u}\,ds
    -
    \frac{1}{u}\int_{t+u}^{t}(t-s)^{\alpha-1}\,ds
  \right]
  \to
  \frac{g_z(t)}{\Gamma(\alpha)}\,t^{\alpha-1}
\end{equation}
as $u\uparrow0$.

For the first integral in \eqref{eq:dq-split-proof-neg}, the same dominated convergence argument as above applies, because for fixed $s\in[0,t)$,
\[
  \frac{(t+u-s)^{\alpha-1}-(t-s)^{\alpha-1}}{u}
  \to
  (\alpha-1)(t-s)^{\alpha-2},
\]
as $u\uparrow0$, and the same integrable domination by
$C_{\eta,R,T}(t-s)^{2\alpha-2}$ holds. Hence,
\begin{align}
  &\int_0^{t+u}
  \frac{(t+u-s)^{\alpha-1}-(t-s)^{\alpha-1}}{u}
  \left(g_z(s)-g_z(t)\right)\,ds
  \nonumber\\
  &\to
  (\alpha-1)\int_0^t (t-s)^{\alpha-2}\left(g_z(s)-g_z(t)\right)\,ds,
  \qquad\text{as }u\uparrow0.
  \label{eq:first-integral-limit-neg}
\end{align}

For the second integral in \eqref{eq:dq-split-proof-neg}, write $v:=-u>0$. Then
\begin{align}
  &\left|
    -\frac{1}{u}\int_{t+u}^{t}(t-s)^{\alpha-1}
    \left(g_z(s)-g_z(t)\right)\,ds
  \right|
  \nonumber\\
  &=
  \frac{1}{v}\left|
    \int_{t-v}^{t}(t-s)^{\alpha-1}
    \left(g_z(s)-g_z(t)\right)\,ds
  \right|
  \nonumber\\
  &\le
  \frac{C_{\eta,R,T}}{v}
  \int_{t-v}^{t}(t-s)^{\alpha-1}(t-s)^\alpha\,ds
 =
  C_{\eta,R,T}v^{2\alpha-1}
  \to0,
\label{eq:second-integral-limit-neg}
\end{align}
as $u\uparrow0$.

Combining the limits in \eqref{eq:bracket-limit-neg}, \eqref{eq:first-integral-limit-neg}, and \eqref{eq:second-integral-limit-neg},
we obtain
\begin{equation}
  \lim_{u\uparrow0}\frac{h(t+u,z)-h(t,z)}{u}
  =
  \frac{g_z(t)}{\Gamma(\alpha)}t^{\alpha-1}
  +
  \frac{\alpha-1}{\Gamma(\alpha)}
  \int_0^t (t-s)^{\alpha-2}\left(g_z(s)-g_z(t)\right)\,ds .
  \label{eq:left-derivative-proof}
\end{equation}

From \eqref{eq:right-derivative-proof} and \eqref{eq:left-derivative-proof}, the right and left derivatives agree. Therefore $h(\cdot,z)$ is differentiable on $(0,T]$, and
\begin{equation}\label{eq:h-derivative-formula-proof}
  \partial_t h(t,z)
  =
  \frac{g_z(t)}{\Gamma(\alpha)}t^{\alpha-1}
  +
  \frac{\alpha-1}{\Gamma(\alpha)}
  \int_0^t (t-s)^{\alpha-2}\left(g_z(s)-g_z(t)\right)\,ds .
\end{equation}

It remains to estimate the derivative. Since $|g_z(t)|\le C_{\eta,R,T}$,
\[
  \left|\frac{g_z(t)}{\Gamma(\alpha)}t^{\alpha-1}\right|
  \le C_{\eta,R,T}t^{\alpha-1}.
\]
For the integral term, using once more the $\alpha$-H\"older continuity of $g_z$,
\begin{align*}
  \left|
    \int_0^t (t-s)^{\alpha-2}\left(g_z(s)-g_z(t)\right)\,ds
  \right|
  \le
  C_{\eta,R,T}\int_0^t (t-s)^{2\alpha-2}\,ds
 =
  C_{\eta,R,T}\frac{t^{2\alpha-1}}{2\alpha-1}.
\end{align*}
Since $t\le T$, we have $t^{2\alpha-1}=t^{\alpha-1}t^\alpha\le T^\alpha t^{\alpha-1}$.
Therefore
\[
  |\partial_t h(t,z)|
  \le C_{\eta,R,T}t^{\alpha-1},
  \qquad 0<t\le T.
\]

Finally, since \eqref{eq:h-derivative-formula-proof} holds for every $t\in(0,T]$
and $|\partial_t h(t,z)|\le C_{\eta,R,T}\,t^{\alpha-1}$ with $t^{\alpha-1}\in L^1(0,T)$ (as $\alpha>0$), we have $\partial_t h(\cdot,z)\in L^1(0,T)$. Hence the representation
\[
  h(t,z)=h(0,z)+\int_0^t \partial_s h(s,z)\,ds,\qquad 0\le t\le T,
\]
which follows by integrating \eqref{eq:h-derivative-formula-proof} from $0$ to $t$ and using the dominated convergence theorem as the lower limit approaches~$0$, implies that $t\mapsto h(t,z)$ is absolutely continuous on $[0,T]$.
\end{proof}

\end{document}
